\documentclass[11pt]{article}
\usepackage[T1]{fontenc}
\usepackage{lmodern,microtype,amsmath,amssymb,amsthm,booktabs}
\microtypesetup{expansion=false}
\usepackage[margin=1.05in,headheight=14pt]{geometry}
\usepackage{xcolor}
\definecolor{linknavy}{RGB}{25,55,125}
\usepackage[colorlinks=true,linkcolor=linknavy,citecolor=linknavy,urlcolor=linknavy]{hyperref}
\usepackage{orcidlink,fancyhdr}
\hypersetup{
  pdftitle={Levinson's method in short intervals and simple zeros of the zeta function},
  pdfauthor={Felipe Santibanez-Leal},
  pdfsubject={A localized Levinson-Conrey detector of distinct sign changes and its consequence for simple critical zeros of the Riemann zeta function in short intervals},
  pdfkeywords={Riemann zeta function, simple zeros, critical line, short intervals, Levinson method, mollifier, interval arithmetic}
}
\newtheorem{theorem}{Theorem}[section]
\newtheorem{proposition}[theorem]{Proposition}
\newtheorem{lemma}[theorem]{Lemma}
\newtheorem{corollary}[theorem]{Corollary}
\theoremstyle{remark}
\newtheorem{remark}[theorem]{Remark}
\newcommand{\R}{\mathbb R}
\newcommand{\C}{\mathbb C}
\newcommand{\Z}{\mathbb Z}
\newcommand{\Vt}{\widetilde V}
\newcommand{\wh}{\widehat w(0)}
\newcommand{\doctype}{Preprint}
\newcommand{\docversion}{v0.02}
\newcommand{\docdate}{4 October 2026}
\newcommand{\versiondoi}{10.5281/zenodo.23132248}
\newcommand{\conceptdoi}{10.5281/zenodo.22984154}
\newcommand{\docrepo}{github.com/fsantibanezleal/CAOS\_RESEARCH}
\newcommand{\headerblock}{%
\begin{center}
\fbox{\parbox{0.94\textwidth}{\small\raggedright
\textbf{\doctype} \quad$\cdot$\quad \docversion, \docdate \quad$\cdot$\quad Licence CC BY 4.0\\[2pt]
CAOS Research open-problems programme \quad$\cdot$\quad problem \texttt{riemann-hypothesis}\\[2pt]
Version DOI \href{https://doi.org/\versiondoi}{\texttt{\versiondoi}}\\[2pt]
Concept DOI (always latest) \href{https://doi.org/\conceptdoi}{\texttt{\conceptdoi}}\\[2pt]
Not peer reviewed. Source code, data and computational records: \href{https://\docrepo}{\texttt{\docrepo}}
}}
\end{center}
\vspace{4pt}}
\pagestyle{fancy}
\fancyhf{}
\fancyhead[L]{\small CAOS Research \doctype}
\fancyhead[R]{\small Levinson's method in short intervals\textperiodcentered\ \docversion}
\fancyfoot[C]{\thepage}
\fancypagestyle{plain}{\fancyhf{}\fancyfoot[C]{\thepage}\renewcommand{\headrulewidth}{0pt}}
\allowdisplaybreaks[1]

\title{Levinson's method in short intervals and\\simple zeros of the zeta function}
\author{Felipe Santib\'a\~nez-Leal\,\orcidlink{0000-0002-0150-3246}\\
\small CAOS Research program\\
\small\texttt{github.com/fsantibanezleal/CAOS\_RESEARCH}}
\date{\docdate}

\begin{document}
\maketitle
\thispagestyle{plain}
\headerblock

\begin{abstract}
For $1/2<\theta<1$, we prove a uniformly shifted mollified second moment
on windows of length $T^\theta$ for mollifier exponents
$\nu<\min\{1/2,(17/33)(2\theta-1)\}$. An exact composite additive-divisor
transformation reduces the signed off-diagonal to inverse-fraction sums.
Their weight depends on one ratio. A charged Mellin multiplier estimate,
independently recovered from exact Hankel symbols and Fourier separation,
allows application of Bettin and Chandee's trilinear theorem. Both residues,
small complex shifts, unbalanced blocks and smooth-window conversion are
retained. The moment supports Conrey's operator polynomial of arbitrary
fixed degree. Localized counting gives a density of distinct sign changes
of Hardy's function, and Wang's pair theorem with a finite parity inequality
converts it to simple critical zeros. An unchanged degree-201 detector,
checked with independent interval arithmetic, gives positive simple-critical
density for every fixed $\theta\in[0.5339,1)$; at $\theta=0.5339$ the lower
density exceeds $0.0003985233159135$. The estimates are asymptotic, give no
effective starting height, use attributed analytic inputs including Wang's
recent preprint, and do not prove the Riemann hypothesis.
\end{abstract}

\section{Introduction}

Let $\rho=\beta+i\gamma$ denote the nontrivial zeros of the Riemann zeta
function. For $T,H>0$ let $N(T,H)$ count the zeros with $T<\gamma\le T+H$,
with multiplicity; let $S(T,H)$ count the simple zeros among them with
$\beta=1/2$; and let $O(T,H)$ count the distinct points $t\in(T,T+H]$ at which
Hardy's function $Z(t)=e^{i\vartheta(t)}\zeta(1/2+it)$ changes sign, that is,
the distinct critical zeros of odd multiplicity. Throughout, $H=T^\theta$ with
a fixed $\theta\in(1/2,1)$ and $L=\log T$.

Selberg proved that $O(T,T^\theta)\gg N(T,T^\theta)$ for every
$\theta>1/2$, and Karatsuba extended this to $\theta>27/82$~\cite{Karatsuba}.
These theorems do not give simple zeros. For simple critical zeros, Steuding
localized Levinson's method with a linear operator polynomial, proved a
short-window mollified moment with error $O(T^{1/3+\epsilon}M^{4/3})$, and
obtained a positive proportion for $\theta\ge0.552$~\cite{Steuding,SteudingThesis};
he also observed that the mollifier range available from the classical
off-diagonal treatment gives only $\theta\ge0.693$ for linear $Q$. Wang's short-interval
pair-correlation theorem~\cite{WangShort} gives the lower density
\begin{equation}\label{eq:wangc}
 c(\theta)=2-\frac\theta2-\frac1{\sqrt2}\cot\frac{\theta}{\sqrt2},
\end{equation}
which is positive for $\theta>\theta_0=0.550193964744\ldots$. The finite
Hilbert--parity inequality of~\cite[Proposition 2.5]{CAOSShort} converts any
lower density $k$ for $O/N$ into the simple-critical density
\begin{equation}\label{eq:h}
 h(\theta;k)=\frac{3+k-\sqrt{(1-k)(9-k-8c(\theta))}}4 .
\end{equation}
With the rank-six Selberg detector of Pearce-Crump~\cite{PearceCrump},
$k_6(\theta)=(\theta-1/2)/(4eC_6)$ with $C_6=0.65663\ldots$, the resulting
term is positive exactly for $\theta>\theta_6$, where
$0.5458837<\theta_6<0.5458838$~\cite[Theorem 1.2]{CAOSShort}. Since the
positivity condition is $k>1-2/(2-c(\theta))$, the onset is governed by the
size of the odd-support density near $\theta=1/2$, where $k_6$ has slope only
$1/(4eC_6)\approx0.140$.

The localized Levinson detector supplies the odd-support density needed
by the parity inequality. The elementary diagonal range is
$\nu<\theta-1/2$. We extend it to
\begin{equation}\label{eq:range}
 0<\nu<\min\{1/2,(17/33)(2\theta-1)\}
\end{equation}
by retaining and estimating the full signed off-diagonal, using the
attributed trilinear theorem of Bettin and Chandee~\cite{BC}.
The labels [D] and [MV] mark mathematical deductions and finite
machine-verified assertions, respectively.

\begin{theorem}[Simple critical zeros from $T^{0.5339}$, D+MV]
\label{thm:main}
For every fixed $\theta\in[0.5339,1)$,
\begin{equation}\label{eq:main}
 \liminf_{T\to\infty}\frac{S(T,T^\theta)}{N(T,T^\theta)}
 \ge\max\left\{0,c(\theta),\frac{c(\theta)+2\kappa_\theta}{3},
 h(\theta;\kappa_\theta)\right\}>0,
\end{equation}
where $\kappa_\theta$ can be any certified detector in
Table~\ref{tab:kappa} whose $\nu$ satisfies \eqref{eq:range}.
In particular the fixed detector with $\nu=349/10000$ has
$\kappa\ge62587441630700832326901419310999470251/(25\cdot10^{38})$ and
\[
 h(5339/10000;\kappa)>0.0003985233159135.
\]
The same fixed detector is admissible for every larger $\theta<1$.
\end{theorem}

The improvement in range is analytical; no new detector optimization
is needed for the displayed consequence. The moment and counting
statements are as follows.

\begin{theorem}[Mollified moment on a short window, D]\label{thm:moment}
Let $\nu$ satisfy \eqref{eq:range}, let $R>0$, and let $P$ be a real polynomial with $P(0)=0$,
$P(1)=1$, and let $Q$ be a real polynomial with $Q(0)=1$. Put
$\sigma_0=1/2-R/L$, $y=T^\nu$, $\Delta=H/L$,
\[
 V(s)=Q\Bigl(-\frac1L\frac{d}{ds}\Bigr)\zeta(s),\qquad
 \psi(s)=\sum_{h\le y}\frac{\mu(h)}{h^{s+1/2-\sigma_0}}
 P\Bigl(\frac{\log(y/h)}{\log y}\Bigr).
\]
Let $w$ be smooth with $0\le w\le1$, $w=1$ on $[T,T+H]$, support in
$[T-\Delta,T+H+\Delta]$ and $w^{(j)}\ll_j\Delta^{-j}$. Then
\[
 \int_\R w(t)\,|V\psi(\sigma_0+it)|^2\,dt=c(P,Q,R,\nu)\,\wh+O(H/L),
\]
\[
 c(P,Q,R,\nu)=1+\frac1\nu\int_0^1\!\!\int_0^1
 \bigl(w_R(v)P'(u)+\nu w_R'(v)P(u)\bigr)^2du\,dv,\qquad w_R(v)=e^{Rv}Q(v).
\]
\end{theorem}

\begin{theorem}[Distinct sign changes, D]\label{thm:signs}
In the setting of Theorem~\ref{thm:moment}, suppose that $Q(x)+Q(1-x)$ is a
nonzero constant. Then
\[
 \liminf_{T\to\infty}\frac{O(T,T^\theta)}{N(T,T^\theta)}\ge
 \kappa(P,Q,R,\nu):=1-\frac1R\log c(P,Q,R,\nu).
\]
\end{theorem}

The constant is Conrey's~\cite{Conrey} in Young's form~\cite{Young}.
The counting argument pays full weight for detector zeros on the critical
line; a polynomial of higher degree supplies distinct odd-order zeros,
which are the appropriate input for the parity inequality. Conrey's
Gaussian-window proposition applies under a different restriction,
$1-\theta<1/7-\nu/4$. Bounded primary-source searches did not locate the
shifted short-window range \eqref{eq:range} or the onset $0.5339$; this
does not establish worldwide priority. The global $17/33$ range is
context~\cite{BC}, not a short-window theorem imported without proof.

\section{The elementary diagonal range}\label{sec:moment}

Throughout this section assume $0<\nu<\theta-1/2$. It supplies the
arithmetic main-term evaluation and localized counting inputs; the
signed estimate of Section~\ref{sec:signed} proves the full range
\eqref{eq:range}.

We follow Young~\cite{Young}. His Theorem 2 treats a weight supported in
$[T/4,2T]$ with $w^{(j)}\ll(T/L)^{-j}$ and a mollifier of length $M=T^{\vartheta}$,
$\vartheta<1/2$. Here the weight has support of length at most $2H$, the
derivative scale is $\Delta=H/L$, and $M$ is replaced by $y=T^\nu$. We list
every step in which the length enters; all others are used verbatim. As in
Young's Section 3, it suffices to prove
\begin{equation}\label{eq:Iab}
 I(\alpha,\beta):=\int w(t)\zeta(\tfrac12+\alpha+it)\zeta(\tfrac12+\beta-it)
 |\psi(\sigma_0+it)|^2dt=c(\alpha,\beta)\,\wh+O(H/L)
\end{equation}
uniformly for $\alpha,\beta\ll1/L$, with $c(\alpha,\beta)$ Young's (3.3)
with $\vartheta$ replaced by $\nu$ and $M$ by $y$. We first work on the annuli
$\alpha,\beta\asymp1/L$, $|\alpha+\beta|\gg1/L$ of Young's Lemma 6. There the
factor $G(s)=e^{s^2}p(s)$, $p(s)=((\alpha+\beta)^2-4s^2)/(\alpha+\beta)^2$, of
Young's approximate functional equation (Lemma 4) satisfies
$p(s)\ll L^2(1+|s|^2)$. Hence the weight $V_{\alpha,\beta}(x,t)$ and the
factor $X_{\alpha,\beta,t}$ of that lemma obey, for $t\ge T/2$,
\begin{equation}\label{eq:Vbound}
 t^j\partial_t^jV_{\alpha,\beta}(x,t)\ll_{A,j}L^2(1+x/t)^{-A},\qquad
 X_{\alpha,\beta,t}=(t/2\pi)^{-\alpha-\beta}(1+O(t^{-1})).
\end{equation}
The first bound is Young's (4.3), whose decay factor is printed as
$(1+|t/x|)^{-A}$; the Mellin representation (4.1) gives $(1+x/t)^{-A}$,
which is the form used in his proof of Lemma 5. The Stirling expansion (4.2)
of $g_{\alpha,\beta}(s,t)$ is uniform only for $|s|=o(t^{1/2})$; on every
vertical line used below, the part $|\Im s|\ge T^{1/3}$ is negligible because
$|G(s)|\ll L^2e^{-(\Im s)^2/2}$ there, while
$|g_{\alpha,\beta}(s,t)|\ll(t+|s|)^{O(1)}$.

\begin{lemma}[Twisted integral]\label{lem:twisted}
Let $h,k\le y$ and let $\alpha,\beta$ lie on the annuli. Then
\begin{multline*}
 \int w(t)\Bigl(\frac hk\Bigr)^{-it}\zeta(\tfrac12+\alpha+it)\zeta(\tfrac12+\beta-it)dt
 =\sum_{hm=kn}\frac{\int V_{\alpha,\beta}(mn,t)w(t)dt}{m^{1/2+\alpha}n^{1/2+\beta}}\\
 +\sum_{hm=kn}\frac{\int V_{-\beta,-\alpha}(mn,t)X_{\alpha,\beta,t}w(t)dt}{m^{1/2-\beta}n^{1/2-\alpha}}
 +O_A(T^{-A}).
\end{multline*}
\end{lemma}

\begin{proof}
Insert the approximate functional equation. Its error term contributes
$O(HT^{-A})$ and the terms with $hm=kn$ are the main terms. For $hm\ne kn$,
the bounds \eqref{eq:Vbound} and $w^{(j)}\ll\Delta^{-j}$ give
$\partial_t^j[w(t)V_{\alpha,\beta}(x,t)]\ll L^2(1+x/T)^{-A}\Delta^{-j}$ on the
support, so $j$ integrations by parts give
\[
 \int w(t)\Bigl(\frac{hm}{kn}\Bigr)^{-it}V_{\alpha,\beta}(mn,t)dt
 \ll HL^2\frac{(1+mn/T)^{-A}}{(\Delta|\log(hm/kn)|)^j}
 \le HL^2(1+mn/T)^{-A}\Bigl(\frac{2\sqrt{hkmn}}{\Delta}\Bigr)^j,
\]
since $|\log(hm/kn)|\ge1/(2\sqrt{hkmn})$ for $hm\ne kn$. Put
$\eta=\theta-1/2-\nu>0$ and $\epsilon=\eta/2$. If $mn\le T^{1+\epsilon}$ then
$2\sqrt{hkmn}/\Delta\le2yT^{(1+\epsilon)/2}L/T^\theta=2LT^{-3\eta/4}$, and there
are at most $T^{2}$ such pairs $(m,n)$, so these terms contribute
$O(T^{-A})$ once $j>4(A+5)/(3\eta)$. If $mn>T^{1+\epsilon}$ we do not integrate
by parts: the trivial bound $HL^2(mn/T)^{-A'}$ summed against
$(mn)^{-1/2+O(1/L)}$ gives $O(HL^3T^{(1+\epsilon)/2-\epsilon A'+o(1)})$, which is
$O(T^{-A})$ for $A'\ge(A+3)/\epsilon$. The reflected part is identical because
$t^j\partial_t^jX_{\alpha,\beta,t}\ll|X_{\alpha,\beta,t}|\ll1$.
\end{proof}

Inserting the mollifier, $I=I_1+I_2+O(T^{-A})$, where $I_1$ is Young's (5.1)
with $M$ replaced by $y$ and $I_2$ is $I_1(-\beta,-\alpha)$ with
$X_{\alpha,\beta,t}$ inserted in the $t$-integral. On the support of $w$ we have
$t=T(1+O(H/T))$, so by \eqref{eq:Vbound}
$X_{\alpha,\beta,t}=(T/2\pi)^{-\alpha-\beta}(1+\rho(t))$ with $\rho\ll H/T$.
Write $h=gh'$, $k=gk'$ with $(h',k')=1$; then $hm=kn$ means $m=k'r$, $n=h'r$,
and
\[
 \sum_{h,k\le y}\frac1{\sqrt{hk}}\sum_{hm=kn}
 \frac{(1+mn/T)^{-A}}{(mn)^{1/2-O(1/L)}}
 \ll L\sum_{g\le y}\frac1g\Bigl(\sum_{h'\le y}\frac1{h'}\Bigr)^2\ll L^4.
\]
Bounding with absolute values, the contribution of $\rho$ to $I_2$ is
$\ll H(H/T)L^6=o(H/L)$ because $\theta<1$. Since $(2\pi)^{\alpha+\beta}=1+O(1/L)$
and $|I_1(-\beta,-\alpha)|\ll H$ by the next lemma (the region of the annuli
is invariant under $(\alpha,\beta)\mapsto(-\beta,-\alpha)$),
\begin{equation}\label{eq:I1I2}
 I(\alpha,\beta)=I_1(\alpha,\beta)+T^{-\alpha-\beta}I_1(-\beta,-\alpha)+O(H/L).
\end{equation}

\begin{lemma}[Diagonal]\label{lem:diagonal}
Uniformly on the annuli, $I_1(\alpha,\beta)=c_1(\alpha,\beta)\wh+O(H/L)$ with
\[
 c_1(\alpha,\beta)=\frac1{(\alpha+\beta)\log y}\int_0^1
 \bigl(P'(u)+\alpha\log y\,P(u)\bigr)\bigl(P'(u)+\beta\log y\,P(u)\bigr)du .
\]
\end{lemma}

\begin{proof}
Follow Young's Section 6. After the Mellin representation (6.1) and the
arithmetic identity (6.2), move the $u,v$-contours to $\Re=\delta$ and the
$s$-contour to $\Re s=-\delta+\epsilon_1$ with a fixed small $\epsilon_1>0$. The only pole crossed is $s=0$,
because $G$ vanishes at the pole $s=-(\alpha+\beta)/2$ of
$\zeta(1+\alpha+\beta+2s)$. The only factor involving $w$ is
$\int w(t)g_{\alpha,\beta}(s,t)dt$, and by Young's (4.2),
$|g_{\alpha,\beta}(s,t)|\ll(t/2\pi)^{-\delta+\epsilon_1}(1+|s|^2)$ on the new
contour for $|\Im s|\le T^{1/3}$. This factor is therefore
$\ll HT^{-\delta+\epsilon_1}(1+|s|^2)$, while $|y^{u+v}|=y^{2\delta}$; the
remaining factors are as in Young. The new contour contributes
$\ll HT^{-\delta(1-2\nu)+\epsilon_1}L^{O(1)}=o(H/L)$ for
$\epsilon_1<\delta(1-2\nu)/2$, since $\nu<1/2$. The residue at $s=0$ has
$g_{\alpha,\beta}(0,t)=G(0)=1$ and equals $\wh$ times the expression in
Young's (6.3). Young's Lemma 7 and Section 7 involve neither $w$ nor $T$ and
give the main term with relative error $O(1/L)$.
\end{proof}

\begin{proof}[Proof of Theorem~\ref{thm:moment}]
As in Young's deduction of his Lemma 3 from Lemma 6,
\[
 c_1(\alpha,\beta)+c_1(-\beta,-\alpha)=\int_0^12PP'=1,\qquad
 \frac{1-T^{-\alpha-\beta}}{(\alpha+\beta)\log y}=\frac1\nu\int_0^1T^{-v(\alpha+\beta)}dv .
\]
With \eqref{eq:I1I2} and Lemma~\ref{lem:diagonal} this gives \eqref{eq:Iab} on
the annuli, and the maximum modulus principle extends it to $\alpha,\beta\ll1/L$
because both sides are holomorphic. By Young's (3.4), the left side of
Theorem~\ref{thm:moment} equals
\[
 Q\Bigl(-\frac1L\frac{\partial}{\partial\alpha}\Bigr)
 Q\Bigl(-\frac1L\frac{\partial}{\partial\beta}\Bigr)I(\alpha,\beta)
 \Big|_{\alpha=\beta=-R/L};
\]
the derivatives are Cauchy integrals over circles of radius $\asymp1/L$ on
which the error is uniform, and the operator applied to $c(\alpha,\beta)$ gives
Young's (1.3) with $\vartheta$ replaced by $\nu$. Since
\[
 \frac{d}{dx}\Bigl[e^{R\nu x}P(x+u)Q(v+\nu x)\Bigr]_{x=0}
 =e^{-Rv}\bigl(w_R(v)P'(u)+\nu w_R'(v)P(u)\bigr),
\]
this is the stated constant.
\end{proof}

\begin{remark}\label{rem:uses}
The hypothesis $\nu<\theta-1/2$ is used only in Lemma~\ref{lem:twisted}, and
the bound $\nu<1/2$ only in Lemma~\ref{lem:diagonal}. The hypothesis
$\theta<1$, through $H=o(T)$, places the support of $w$ in $[T/2,2T]$, so that
$t^{-1}\le\Delta^{-1}$, and bounds $\rho$. The normalization $Q(0)=1$ enters
only through the constant term: the operator maps the term $1$ of
$c(\alpha,\beta)$ to $Q(0)^2$, so for every real polynomial $Q$ the moment is
$\bigl(Q(0)^2+\nu^{-1}\int_0^1\!\int_0^1(w_RP'+\nu w_R'P)^2\bigr)\wh+O(H/L)$,
and in particular $O(H)$.
\end{remark}


\section{The signed estimate for a longer mollifier}\label{sec:signed}

All orders in this section are fixed before $T$ tends to infinity.
Write $e(x)=e^{2\pi ix}$ and
$\sigma_{a,b}(n)=\sum_{uv=n}u^{-a}v^{-b}$. Shifts lie in a fixed
$C/\log T$ disk, enlarged by a fixed factor when taking derivatives.

\subsection{An exact Gaussian reduction and its summable remainder}
For the Gaussian of width $H$, define
\[
 K_\beta(x)=\frac{2x^{-\beta}}{\sqrt\pi}
 \int_\R e^{-y^2+(1-2\beta+2iT)y/H}
             \cos(2\pi x e^{2y/H})\,dy.
\]
Functional equation and Mellin inversion give
\begin{align}\label{eq:exactGaussian}
 &\int_\R\zeta(\tfrac12+\alpha+it)\zeta(\tfrac12+\beta-it)
      (h/k)^{it}e^{-(t-T)^2/H^2}\,dt\notag\\
 &=2\pi\sqrt{k/h}\sum_{n\ge1}\sigma_{\alpha,-\beta}(n)K_\beta(nk/h)
   -2\pi\zeta(\alpha+\beta)G(1-\alpha)(h/k)^{1/2-\alpha},
\end{align}
where $G(s)=\exp((s-1/2-iT)^2/H^2)$. For justification substitute
$u=e^{2y/H}$ in the kernel and regularize its $x$ Mellin transform by
$e^{-\epsilon x}$. Dominated convergence and Gaussian integration give
\[
 \widetilde K_\beta(s)=2(2\pi)^{-(s-\beta)}
       \Gamma(s-\beta)\cos(\pi(s-\beta)/2)G(s).
\]
Move the moment's $s$ contour from $1/2$ to the absolute Dirichlet
half-plane. The pole $s=1-\alpha$ is crossed with the displayed minus
sign. The apparent pole $1+\beta$ is canceled by the cosine zero.
The kernel has arbitrary inverse-power decay at infinity; its defining
log-Gaussian amplitude and all derivatives have integrable endpoints.

On a fixed central band $x\asymp T$, expand only the nonlinear phase
$\exp[\epsilon2\pi ix(e^{2y/H}-1-2y/H)]$ to order $J-1$.
Keep the linear Fourier phase exact. The $j$th profile is
\[
 A_{\epsilon,j}(x)=\frac{x^{-\beta}(\epsilon2\pi ix)^j}{j!\sqrt\pi}
 \int_\R e^{-y^2+(1-2\beta)y/H}
 e^{2i(T+\epsilon2\pi x)y/H}(e^{2y/H}-1-2y/H)^j\,dy.
\]
Multiply by a smooth central cutoff. Taylor's real-phase remainder is
$O_J(x^{-\Re\beta}(x/H^2)^J)$. Outside the band, nonstationary
integration by parts after the $u$ substitution bounds the kernel by
$O_A(x^{-\Re\beta}(H/(T+x))^A)$, plus an exponentially small summable
tail. In a bounded mollifier average the resulting absolute error is
\begin{equation}\label{eq:representationerror}
 O_{J,A,\epsilon}\left(T^{1+\nu+(1+\nu)\epsilon}\log T
       \{(T/H^2)^J+(H/T)^A\}\right).
\end{equation}
Indeed the central count per $(h,k)$ is $O(Th/k)$, the outer factor is
$1/h$, and $\sigma_{\alpha,-\beta}(n)\ll_\epsilon n^\epsilon$ there.
The absolute $h,k$ sum is $O(M\log M)$ after taking out $T$.
For the tail sum integrate $(T+(k/h)n)^{-A}$, keeping its scale
$Th/k\ge T/M$. Thus fixed large $J,A$ make
\eqref{eq:representationerror} smaller than $H/\log^B T$ for every
prescribed fixed $B$. This uses no signed cancellation.
The retained kernel is
$\sum_{\epsilon=\pm1}e(\epsilon x)\sum_{j<J}A_{\epsilon,j}(x)$.
Thus each divisor sum has the additive phase $e(\epsilon np/q)$;
the smooth amplitudes above do not include that arithmetic phase.

\subsection{Composite transformation, with both residues}
Put $h=gq,k=gp$, $(p,q)=1$. The actual outer coefficient is
$b_{gq}\overline b_{gp}/(gq)$, where $b_m=\mu(m)P(\log(M/m)/\log M)$.
Squarefree support retains $(g,pq)=1$; neither weight nor cutoff is removed.
Apply the following transformation with additive numerator $A=\epsilon p$
and smooth weight $F(n)=A_{\epsilon,j}(np/q)$ for each retained profile.
For $(A,q)=1$ set
$D_{a,b}(s,A/q)=\sum_n\sigma_{a,b}(n)e(An/q)n^{-s}$.
Opening both residue variables gives
\[
 D_{a,b}(s,A/q)=q^{-2s-a-b}
 \sum_{r,t=1}^q e(Art/q)\zeta(s+a,r/q)\zeta(s+b,t/q).
\]
The Hurwitz Fourier formula~\cite{DLMF} and the exact full-residue sum
\[
 \sum_{r,t\bmod q}e((Art+\epsilon nr+\eta mt)/q)
       =q\,e(-\epsilon\eta nm\overline A/q)
\]
give, with $z=a+b,d=a-b$,
\begin{align}\label{eq:Estermann}
 D_{a,b}(s,A/q)&=2q^{1-2s-z}(2\pi)^{2s+z-2}
 \Gamma(1-s-a)\Gamma(1-s-b)\notag\\
 &\quad\cdot[\cos(\pi d/2)D_{-a,-b}(1-s,\overline A/q)
 -\cos(\pi(s+z/2))D_{-a,-b}(1-s,-\overline A/q)].
\end{align}
Sum $r$ first to prove the finite identity; exactly one $t$ survives,
including nonunit $n$. Both Dirichlet series initially converge absolutely
in the appropriate left half-plane, and continuation is meromorphic.
For compact smooth $F$, moving the Mellin contour crosses precisely
\[
 q^{-1+a-b}\zeta(1-a+b)\widetilde F(1-a)
 +q^{-1-a+b}\zeta(1+a-b)\widetilde F(1-b).
\]
At $a=b$ combine the poles before taking the limit: the residue is
$q^{-1}\int F(x)x^{-a}(\log x+2\gamma-2\log q)\,dx$.
A contour enclosing both poles is holomorphic in the shifts.

\subsection{The absolute Mellin multiplier}
Take $a=\alpha,b=-\beta$, now $z=\alpha-\beta,d=\alpha+\beta$,
and $F(n)=A(np/q)$ for any of the finitely many profiles.
Its Gaussian Fourier definition implies, for fixed $\sigma>0$ and $B$,
\begin{equation}\label{eq:profileMellin}
 |\widetilde A(1-\sigma-it-z/2)|
 \ll_{B,\sigma} H T^{-\sigma}(1+|t|H/T)^{-B}.
\end{equation}
To see the uniformity, the Gaussian times
$(e^{2y/H}-1-2y/H)^j$ has every fixed Schwartz seminorm
$O_j(H^{-2j})$. Its Fourier transform is uniformly Schwartz on the
scale $(T+\epsilon2\pi x)/H$. The $x^j$ prefactor contributes
$(T/H^2)^j\le1$. Every logarithmic $x$ derivative costs at most $T/H$;
Mellin integration by parts gives \eqref{eq:profileMellin}.
The plus-frequency primal branch is power-small with arbitrary order.

Change $u=1-s-z/2$ in \eqref{eq:Estermann}. Including $2\pi/(gq)$,
the oscillatory dual profile is exactly
\begin{equation}\label{eq:dualMellin}
 \frac{4\pi}{gpq}n^{-z/2}(q/p)^{-z/2}
 \int_{(\sigma)}\Gamma(u-d/2)\Gamma(u+d/2)\cos(\pi u)
  \left(\frac{pq}{4\pi^2n}\right)^u
  \widetilde A(1-u-z/2)\frac{du}{2\pi i}.
\end{equation}
Start where the dual series converges absolutely. For each finite dyadic
block move to any fixed $\sigma>0$; the gamma poles lie at
$\Re u\le|\Re d|/2$. No pole is crossed for large $T$.
The eventual summable tail bounds justify recombining all blocks.
Uniform Stirling bounds give
\[
 |\Gamma(\sigma+it-d/2)\Gamma(\sigma+it+d/2)\cos(\pi(\sigma+it))|
       \ll_\sigma(1+|t|)^{2\sigma-1}.
\]
Consequently the absolute multiplier norm in \eqref{eq:dualMellin} is
\begin{equation}\label{eq:multipliernorm}
 \ll H T^{-\sigma}(T/H)^{2\sigma}.
\end{equation}
The factor $(pq/n)^{it}$ separates into three unimodular sequences.
Apart from harmless small-shift powers the scale is therefore
$(H/g)(pq)^{\sigma-1}(T/H^2)^\sigma n^{-\sigma}$.
For an unbounded dual $n$ range, charge $n^{C/\log T}\le n^\epsilon$;
it is not assumed bounded independently of $n$.

\subsection{All dyadic blocks and the attributed trilinear estimate}
On $p\asymp P,q\asymp Q,n\asymp N$ the coefficient norms are
$O(P^{\sigma-1/2}),O(Q^{\sigma-1/2}),O_\epsilon(N^{1/2-\sigma+\epsilon})$.
Bettin--Chandee Theorem 1~\cite{BC}, with phase parameter $\pm1$, gives
the product of these norms times
\[
 (1+N/(PQ))^{1/2}
 \{(NPQ)^{7/20+\epsilon}(P+Q)^{1/4}
 +(NPQ)^{3/8+\epsilon}(N(P+Q))^{1/8}\}.
\]
Its arbitrary complex coefficient scope includes the separated twists;
its final reciprocity argument covers both relative orders of $P,Q$.
At $\sigma=1/4$, the ratio to $H$ is bounded by
\begin{align*}
 \frac1g(T/H^2)^{1/4}
 \{N^{3/5}(PQ)^{1/10}(P+Q)^{1/4}
       +N^{3/4}(PQ)^{1/8}(P+Q)^{1/8}\}
 (1+N/(PQ))^{1/2}T^\epsilon.
\end{align*}
Put $N_0=TPQ/H^2$. Below $N_0$, both $N$ powers are positive and
$N_0/(PQ)<1$. Above it move to $\sigma=1/4+A$, gaining $(N_0/N)^A$.
A fixed large $A$ absorbs both powers, divisor losses and the growing
source parameter. If $N_0<1$ every integer block is in the tail and gives
a stronger bound. Thus all dual blocks contribute
\[
 \ll\frac Hg T^\epsilon
 \{(T/H^2)^{17/20}(PQ)^{7/10}(P+Q)^{1/4}
       +(T/H^2)(PQ)^{7/8}(P+Q)^{1/8}\}.
\]
Use $P,Q\le M/g$, without assuming balance. The $g$ sums have powers
$-1-33/20$ and $-1-15/8$ and converge. Logarithmic decompositions are
absorbed in $\epsilon$, giving
\begin{equation}\label{eq:newerror}
 O_\epsilon\left(HT^\epsilon
 \{T^{(17/20)(1-2\theta)+(33/20)\nu}
                  +T^{1-2\theta+(15/8)\nu}\}\right).
\end{equation}
At $p=1$ or $q=1$, a convention for inversion modulo one is unnecessary:
$N_0\le T^{1+\nu-2\theta}<1$ in \eqref{eq:range}; a sufficiently large
contour gives an arbitrarily strong saving by absolute summation.
The nonoscillatory branch has $\cos(\pi d/2)$ instead of $\cos(\pi u)$.
Its gamma product retains exponential vertical decay. Absolute summation
on large fixed $\sigma$ gives $O(H(M^2/T)^\sigma)$, negligible since
$\nu<1/2$. The small plus-frequency primal profile is handled with the
same bounds and an arbitrarily large Schwartz order.

\subsection{Exact Hankel symbols: an independent norm derivation}
For $x\ge1$ and $d$ in a compact subset of $|\Re d|<1/2$, define
\[
 a_\pm(x,d)=\frac1{\Gamma(d+1/2)}\int_0^\infty
 e^{-t}t^{d-1/2}(1\pm it/(2x))^{d-1/2}\,dt.
\]
The $U$ integral and its $K$ identity~\cite[13.4.4, 13.6.10]{DLMF},
continued to $z=\mp ix$, and the connection formula 10.27.8 give exactly
\[
 H_d^{(1,2)}(x)=\sqrt{2/(\pi x)}
 e^{\pm i(x-\pi d/2-\pi/4)}a_\pm(x,d).
\]
For every fixed $k$, $x^k\partial_x^ka_\pm=O_k(1)$ uniformly.
After differentiation the ratio
$(t/x)/|1\pm it/(2x)|\le2$ bounds every extra power; the remaining
integrand is dominated by $e^{-t}t^{\Re d-1/2}$ uniformly. Thus no
unaccounted asymptotic remainder or singular order quotient occurs.
The oscillatory Bessel combination from \eqref{eq:dualMellin} is
\[
 \sin(\pi d/2)J_d(x)+\cos(\pi d/2)Y_d(x)
 =\sqrt{2/(\pi x)}\left[-\tfrac i2e^{i(x-\pi/4)}a_+
                        +\tfrac i2e^{-i(x-\pi/4)}a_-\right].
\]
After $x=T/(2\pi)+Hv/(2\pi)$, the normalized profile depends only on
$r=n/(pq)$. Set $\lambda=\sqrt{NT/(PQ)}$, $\mu=H\sqrt{N/(TPQ)}$.
The phase $4\pi\sqrt{rx}$ has $v$ derivative comparable to $\mu$;
its higher derivatives divided by the first are bounded. Gaussian
Schwartz bounds and integration by parts give, on fixed log-r support,
\[
 \|W\|_2+(1+\lambda)^{-1}\|W'\|_2\ll_A(1+\mu)^{-A}.
\]
With $L_1=1+\lambda$, Fourier Cauchy--Schwarz and Parseval give
\[
 \frac1{2\pi}\int|\widehat W(t)|\,dt
 \le\sqrt{L_1/2}(\|W\|_2^2+L_1^{-2}\|W'\|_2^2)^{1/2}.
\]
Including the actual normalization
$H/[g(PQ)^{3/4}T^{1/4}]$ and the $n^{-1/4}$ coefficient, the three
norms are $O(P^{1/2}),O(Q^{1/2}),O_\epsilon(N^{1/4+\epsilon})$.
The resulting bound before summing $N$ is
\[
 \begin{aligned}
 &\frac Hg\{N^{17/20}(PQ)^{-3/20}(P+Q)^{1/4}
      +N(PQ)^{-1/8}(P+Q)^{1/8}\}\\
 &\hspace{25pt}\cdot(1+N/(PQ))^{1/2}(1+\sqrt{N/N_0})^{-A}T^\epsilon.
 \end{aligned}
\]
Its dominant scale is again $N_0$, reproducing \eqref{eq:newerror}.
This route is sharper below $N_0$; both routes agree at the scale that
determines the final exponents.

\subsection{Residues, compact windows and general operator polynomials}
The leading negative-frequency Gaussian has, after
$x=T/(2\pi)+Hv/(2\pi)$, amplitude
$\exp(-v^2-i(1-2\beta)v/H+(1-2\beta)^2/(4H^2))$.
Its whole-v integral is exactly $\sqrt\pi$.
The two residue terms, including the outer factor, are consequently
\[
 \begin{aligned}
 \frac{H\sqrt\pi}g\big[&
 \zeta(1+\alpha+\beta)q^{-1-\alpha}p^{-1-\beta}\\
 &+(T/(2\pi))^{-\alpha-\beta}\zeta(1-\alpha-\beta)
                       q^{-1+\beta}p^{-1+\alpha}\big]
 \end{aligned}
\]
up to power-small errors. Replacing the $x$ power by its centered value
costs $O(H/T)$ times logarithms; $j\ge1$ profiles cost $(T/H^2)^j$.
The absolute outer sum has only logarithmic size because its weights
are $1/(gpq)$. At $\alpha+\beta=0$, combine the residues; maximum-modulus
bounds on a circle of radius $C'/\log T$ avoid a false pole loss.
Fixed shift derivatives cost fixed logarithmic powers. Young's arithmetic
evaluation (Lemma~\ref{lem:diagonal} and his Lemma 7) supplies the same
Conrey constant for $\nu<1/2$. In particular the residue sum is the full
diagonal main term, not an additional conjectured signed term.

For compact windows choose fixed $\eta>0$ so that
$\theta-\eta>1/2$ and $\nu<(17/33)(2(\theta-\eta)-1)$.
Let $s_0=T^{\theta-\eta}$ and
$\phi_{s_0}(t)=e^{-t^2/s_0^2}/(\sqrt\pi s_0)$. For fixed large $D$ set
\[
 w_D=\phi_{s_0}*\sum_{j<D}\frac{(-s_0^2/4)^j}{j!}w^{(2j)}.
\]
The Fourier deficit $1-e^{-u}\sum_{j<D}u^j/j!$ lies between $0$ and
$u^D/D!$. Thus $\|w-w_D\|_\infty\ll_D
T^{-2D\eta}\log^{2D+O(1)}T$ for the stated logarithmic edge derivatives.
A coarse polynomial zeta-mollifier bound on $|t|\le2T$ and Gaussian
tails elsewhere give a power-small replacement error for sufficiently
large fixed $D$, without using the desired moment. The signed mixture
has absolute coefficient norm
\[
 O\left((H/s_0)\sum_{j<D}(s_0\log T/H)^{2j}\right)=O(H/s_0).
\]
Gaussian errors $O(s_0T^{-\delta})$ therefore become $O(HT^{-\delta})$.
The true loss is evaluating \eqref{eq:newerror} at $\theta-\eta$;
it cannot be omitted. Window mass is preserved exactly because all
nonzero derivatives integrate to zero. Centers vary by $O(H)$;
replacing their logarithmic normalization by $\log T$ costs only
power-small errors. Applying the uniform shifted asymptotic by Cauchy
to every fixed-degree $Q$ completes Theorem~\ref{thm:moment} in
\eqref{eq:range}. Neither a sharp indicator nor an infinite Taylor
expansion has been substituted without an error estimate.

\section{Distinct sign changes}\label{sec:signs}

\subsection{An exact identity on the critical line}

Let $\chi(s)=\pi^{s-1/2}\Gamma((1-s)/2)/\Gamma(s/2)$, so that
$\zeta(s)=\chi(s)\zeta(1-s)$, and put $\lambda(s)=-\chi'(s)/\chi(s)$. Write
$D=d/ds$ and $Q(x)=\sum_kq_kx^k$.

\begin{lemma}\label{lem:identity}
For every $k\ge0$, $\chi(s)\zeta^{(k)}(1-s)=(-1)^k(D+\lambda)^k\zeta(s)$, where
$(D+\lambda)^k$ is the $k$-fold composition of $f\mapsto f'+\lambda f$.
Consequently $\chi(s)V(1-s)=\sum_kq_kL^{-k}(D+\lambda)^k\zeta(s)$. If
$Q(x)+Q(1-x)=\beta$ identically, then
\begin{equation}\label{eq:Ebeta}
 E_\beta:=\beta\zeta-V-\chi(s)V(1-s)=\sum_kq_kL^{-k}
 \bigl[(D+L)^k-(D+\lambda)^k\bigr]\zeta(s),
\end{equation}
and with $\Vt=V+E_\beta/2$ and $\omega(t)=e^{i\vartheta(t)}$,
\begin{equation}\label{eq:real}
 \beta Z(t)=2\,\Re\bigl(\omega(t)\Vt(\tfrac12+it)\bigr)\qquad(t\in\R).
\end{equation}
\end{lemma}

\begin{proof}
Put $f(s)=\zeta(1-s)=\chi(1-s)\zeta(s)$ and note $-\chi'(1-s)/\chi(1-s)=
\lambda(s)$ because $\chi(s)\chi(1-s)=1$. Then $f'=\chi(1-s)(D+\lambda)\zeta$ and,
inductively, $f^{(k)}=\chi(1-s)(D+\lambda)^k\zeta$; since
$\zeta^{(k)}(1-s)=(-1)^kf^{(k)}(s)$, the first claim follows, and the second by
linearity. The polynomial identity $Q(-x)+Q(1+x)=\beta$ with $x=D/L$ gives
$\beta\zeta=V+\sum_kq_kL^{-k}(D+L)^k\zeta$, which is \eqref{eq:Ebeta}. On
$s=1/2+it$ we have $1-s=\bar s$, $V(\bar s)=\overline{V(s)}$ and
$\chi(s)=\omega^{-2}$, so $\omega\chi(s)V(1-s)=\overline{\omega V(s)}$ and
$\omega\zeta=Z$. Hence $\beta Z=2\Re(\omega V)+\omega E_\beta$. The left side
and $2\Re(\omega V)$ are real, so $\omega E_\beta$ is real, and
\eqref{eq:real} follows.
\end{proof}

The identity \eqref{eq:real} holds for every real $Q$; the symmetry of $Q$
makes $E_\beta$ small. By Stirling's formula, $\lambda(s)=\log(|t|/2\pi)+
O(1/|t|)$ and $\lambda^{(j)}(s)\ll1/|t|$ for $j\ge1$, uniformly in fixed
vertical strips, so $\lambda-L=O(1)$ for $t\asymp T$. In normal-ordered form
$(D+\lambda)^k=\sum_{m\le k}c_{k,m}D^m$, where $c_{k,m}$ is
$\binom km\lambda^{k-m}$ plus a polynomial in $\lambda,\lambda',\dots$ of
weight $k-m$ whose terms contain a derivative of $\lambda$. Hence the
coefficient of $D^m$ in $(D+L)^k-(D+\lambda)^k$ is $O(L^{k-m-1})$, and
\begin{equation}\label{eq:Esize}
 E_\beta=\sum_{j<\deg Q}\varepsilon_j(s)L^{-j}\zeta^{(j)}(s),\qquad
 \varepsilon_j(s)\ll1/L,
\end{equation}
uniformly for $\sigma_0\le\sigma\le\sigma_1$ and $T/2\le t\le2T$. For $j\ge1$,
$L^{-j}\zeta^{(j)}=(-1)^j(V_{1+x^j}-\zeta)$, where $V_{1+x^j}$ is the $V$ of
the polynomial $1+x^j$, so Theorem~\ref{thm:moment} for the polynomials $1$
and $1+x^j$ and Cauchy--Schwarz give
\begin{equation}\label{eq:Emean}
 \int w\,|\psi E_\beta(\sigma_0+it)|^2dt\ll H/L^2 .
\end{equation}
Published admissible polynomials often have $\beta\ne1$; for instance
Conrey's polynomial of degree seven has $\beta=0.984$. For them the
unnormalized difference $\zeta-V-\chi V(1-s)$ is not small, which is why
$E_\beta$ is normalized by $\beta$.

\subsection{The counting lemma}

\begin{lemma}\label{lem:count}
Let $\sigma_1=\sigma_1(P,Q)$ be fixed and large, let $L=\log T_0$ be fixed with
$T_0$ large, and let $|T-T_0|\le1$, $H\asymp T_0^\theta$, where $T$ and $T+H$
are not ordinates of zeros of $Z$ or of $\Vt$ in $[1/2,\sigma_1]$. Let
$N_{\Vt}$ count, with multiplicity, the zeros of $\Vt$ in
$[1/2,\sigma_1]\times(T,T+H]$, including those on $\sigma=1/2$. Then
\[
 O(T,H)\ge N(T,H)-2N_{\Vt}-O(L).
\]
\end{lemma}

\begin{proof}
Let $t_1<\dots<t_J$ be the zeros of $\Vt(1/2+it)$ in $(T,T+H)$, of orders
$k_1,\dots,k_J$, and $K=\sum k_j$. Write
$\Vt(s)=\prod_j(s-1/2-it_j)^{k_j}F(s)$ with $F$ holomorphic and zero-free on the
segment. On the line $\omega\Vt=i^Kp(t)\omega F$ with
$p(t)=\prod_j(t-t_j)^{k_j}$ real, so by \eqref{eq:real}
$\beta Z=2p\,\Re W_0$, where $W_0=i^K\omega F$ is continuous and zero-free on
$[T,T+H]$.

Let $\phi$ be a continuous argument of $W_0$; assume $\phi(T)<\phi(T+H)$, the
other case being symmetric, and let $c_1<\dots<c_M$ be the points of
$\pi/2+\pi\Z$ in $(\phi(T),\phi(T+H))$, so that $M\ge|\Delta\phi|/\pi-1$. Put
$\tau_0=T$, $\tau_M=T+H$, and for $0<i<M$ let $\tau_i$ be the first point with
$\phi(\tau_i)=c_i+\pi/2$. Then $\tau_0<\dots<\tau_M$ and $\cos\phi(\tau_i)$ alternates in
sign, since $\phi(\tau_0)\in(c_1-\pi,c_1)$, $\phi(\tau_M)\in(c_M,c_M+\pi)$ and
$\phi(\tau_i)=c_i+\pi/2$. After moving the interior $\tau_i$ slightly, none of them
is a $t_j$. As $\operatorname{sgn}Z(\tau_i)=\operatorname{sgn}(\beta)
\operatorname{sgn}p(\tau_i)\operatorname{sgn}\cos\phi(\tau_i)$ and $p$ changes sign
only at the $t_j$, the function $Z$ has opposite signs at the ends of at least
$M-J$ of the disjoint intervals $(\tau_{i-1},\tau_i)$, and each such interval
contains a sign change. Hence $O(T,H)\ge|\Delta\phi|/\pi-1-J$.

Apply the argument principle to $F$ on $[1/2,\sigma_1]\times[T,T+H]$. It is
zero-free on the left side, and its zeros inside are the $N_>$ zeros of $\Vt$
with $\sigma>1/2$. On the bottom, right and top sides
$\arg F=\arg\Vt-\sum_jk_j\arg(s-1/2-it_j)$, and each factor turns by exactly
$+\pi$ along these sides. Therefore
\[
 \Delta\phi=\Delta\vartheta+\Delta_{\rm brt}\arg\Vt-\pi K-2\pi N_> .
\]
The Riemann--von Mangoldt formula gives $\Delta\vartheta/\pi=N(T,H)+O(L)$. On
$\sigma=\sigma_1$ we have $|V-1|\le1/4$ and $|E_\beta|\ll1/L$, so $\Re\Vt>0$ and
that side contributes $O(1)$. On the horizontal sides the Backlund--Jensen
argument gives $O(L)$, because $\Vt=(\beta\zeta+V-\chi V(1-s))/2=
\sum_{j\le\deg Q}a_j(s)\zeta^{(j)}(s)$ with each $a_j$ a polynomial in
$L^{-1}$, $\lambda$ and its derivatives, so $\Vt$ is holomorphic and
polynomially bounded near the window and $\Re\Vt\ge1/2$ at $\sigma_1+iT$.
Since $J\le K$ and $N_{\Vt}=N_>+K$,
\[
 O(T,H)\ge N(T,H)-2(N_>+K)-O(L)=N(T,H)-2N_{\Vt}-O(L).\qedhere
\]
\end{proof}

For a general window we keep $L=\log T$ and move both endpoints inward by less
than one unit to avoid the finitely many exceptional ordinates, including those
of zeros of $\psi\Vt$; each unit interval contains $O(L)$ zeros of $\zeta$ and,
by Jensen's formula, $O(L)$ zeros of $\Vt$ and of $\psi\Vt$ in
$[\sigma_0,\sigma_1]$, so every count changes by $O(L)$. The quantity $J$ in
the proof may be replaced by the number of $t_j$ of odd order.

A double zero of $Z$ is not a sign change and is never counted; it is paid
for in $N_{\Vt}$, as an on-line zero of $\Vt$ when $\deg Q=1$ and, in the
examples of Section~\ref{sec:certify}, as a zero of $\Vt$ to the right of the
line when $\deg Q=3$. If the on-line zeros of $\Vt$ are omitted from
$N_{\Vt}$, the inequality fails (Section~\ref{sec:certify}).

\subsection{Proof of Theorem~\ref{thm:signs}}

Let $F_1=\psi\Vt$. Littlewood's lemma on $[\sigma_0,\sigma_1]\times[T,T+H]$
gives
\[
 2\pi\int_{\sigma_0}^{\sigma_1}n(\sigma)d\sigma=
 \int_T^{T+H}\log|F_1(\sigma_0+it)|dt-\int_T^{T+H}\log|F_1(\sigma_1+it)|dt+O(L),
\]
where $n(\sigma)$ counts the zeros of $F_1$ in the rectangle with real part
greater than $\sigma$, and $O(L)$ bounds the horizontal argument integrals:
$\psi$ and $\Vt$ are polynomially bounded and have real part at least $1/2$ at
$\sigma_1+iT$ and $\sigma_1+i(T+H)$, so the Backlund--Jensen argument bounds
the variation of $\arg\psi$ and of $\arg\Vt$, hence of $\arg F_1$. Each zero counted in $N_{\Vt}$ has real
part at least $1/2=\sigma_0+R/L$, and every other zero contributes
nonnegatively, so the left side is at least $2\pi(R/L)N_{\Vt}$.

On $\sigma=\sigma_1$, $\psi V=1+\sum_{n\ge2}b_nn^{-s}$ with
$|b_n|\ll d(n)(\log n)^{\deg Q}$, so its logarithm is an absolutely
convergent Dirichlet series for $\sigma_1$ large and
$\int\log|\psi V(\sigma_1+it)|dt=O(1)$; moreover $|E_\beta/(2V)|\ll1/L$ there,
which contributes $O(H/L)$. On $\sigma=\sigma_0$, concavity of the logarithm,
$w=1$ on $[T,T+H]$ with $w\ge0$, the inequality
$|a+b|^2\le(1+L^{-1/2})|a|^2+(1+L^{1/2})|b|^2$, Theorem~\ref{thm:moment} and
\eqref{eq:Emean} give
\[
 \int_T^{T+H}\log|F_1(\sigma_0+it)|dt\le\frac H2\log\Bigl(\frac1H\int
 w|\psi\Vt|^2\Bigr)\le\frac H2\log c(P,Q,R,\nu)+O(HL^{-1/2}).
\]
Hence $N_{\Vt}\le(HL/4\pi R)\log c+o(HL)$. Since $N(T,H)=(HL/2\pi)(1+o(1))$,
Lemma~\ref{lem:count} gives $O/N\ge1-R^{-1}\log c-o(1)$. \qed

\begin{remark}
The number $\sigma_1$ and all implied constants depend on $P$, $Q$, $R$ and
$\nu$. For the polynomials of Section~\ref{sec:certify} the monomial
coefficients of $Q$ have size about $10^{120}$ and $R$ is about $28$, so the
constants are very large, and for $\eta=\theta-1/2-\nu=10^{-4}$ the factor
$2LT^{-3\eta/4}$ of Lemma~\ref{lem:twisted} is below one only when $\log T$
exceeds about $1.7\cdot10^5$. The results are purely asymptotic.
\end{remark}

\section{Transfer to simple zeros}\label{sec:transfer}

\begin{theorem}[D]\label{thm:transfer}
If $\liminf O(T,T^\theta)/N(T,T^\theta)\ge k\ge0$, then
$\liminf S(T,T^\theta)/N(T,T^\theta)\ge h(\theta;k)$.
\end{theorem}

\begin{proof}
This is the argument of~\cite[proof of Theorem 1.2]{CAOSShort} with the
odd-support input replaced by $k$. For a real even $\eta\in
C_c^\infty((-\lambda/2,\lambda/2))$, $0<\lambda<\theta$, with $\int\eta^2=1$,
let $Q_T$ be the pair sum of $f=\eta^2$ over the scaled zeros of the window;
it is unrelated to the polynomial $Q$ of the previous sections. With
$q_T,s_T,o_T$ the ratios of $Q_T,S,O$ to $N$, the Hilbert--parity inequality
\cite[Proposition 2.5]{CAOSShort} reads $(q_T-s_T)(1-o_T)\ge2(1-s_T)^2$, and
Wang's theorem~\cite[Theorem 2.2]{WangShort} gives $q_T\to\mathcal C(f)$. If
$O=N$, every zero of the window is simple and critical and $s_T=1$. Otherwise
$1-o_T>0$ and the inequality gives $q_T\ge s_T$; for every $\epsilon>0$ and all
large $T$, $q_T\le\mathcal C(f)+\epsilon$ and $o_T\ge k-\epsilon$, and enlarging
the nonnegative factors gives
$2(1-s_T)^2\le(\mathcal C(f)+\epsilon-s_T)(1-k+\epsilon)$. Hence
$\liminf s_T\ge R(\mathcal C(f),1-k)$ with
$R(A,b)=(4-b-\sqrt{b(b+8(A-1))})/4$, the smaller root. Letting $f$ tend to the
cosine density at fixed $\lambda$, so that $\mathcal C(f)\to2-c(\lambda)$, and
then $\lambda\to\theta$, gives $R(2-c(\theta),1-k)=h(\theta;k)$.
\end{proof}

For fixed $k<1$, $h(\theta;k)$ increases with $c$ and with $k$, since
$\partial_k[(1-k)(9-k-8c)]=-(10-2k-8c)<0$ for $c<1$. As
$c'(\theta)=\frac12\cot^2(\theta/\sqrt2)>0$, if $h(\theta_*;\kappa)>0$ for a
detector satisfying \eqref{eq:range} at $\theta_*$, then the same detector is admissible and
$h(\theta;\kappa)>0$ for every $\theta\in[\theta_*,1)$. The linear branch
$(c+2k)/3$ in \eqref{eq:main} is the transfer of~\cite[Theorem 1.4]{CAOSShort}
with the same input.

\section{Certified detector constants}\label{sec:certify}

We use $Q$ in the Chebyshev family
\[
 Q(y)=1-y+\sum_{j=1}^{100}x_j\bigl(T_{2j+1}(1-2y)-(1-2y)\bigr),
\]
which has degree $201$, $Q(0)=1$, $Q(1)=0$ and $Q(y)+Q(1-y)=1$ for all rational
$x_j$, because $T_{2j+1}$ is odd with $T_{2j+1}(1)=1$,
and $P(x)=S(rx)/S(r)$ with $S$ the Taylor polynomial of $\sinh$ of degree
$13$. The coefficients $x_j$ are rationals with denominator $10^{40}$ obtained
by minimizing $c$ numerically, and $R$, $r$, $\nu$ are rational. With
$B_P=\int_0^1P^2$, $C_P=\int_0^1P'^2$ and $\int_0^1PP'=1/2$, and since
$\int_0^1w_Rw_R'=(e^{2R}Q(1)^2-Q(0)^2)/2=-1/2$,
\[
 c(P,Q,R,\nu)=\frac12+\frac{C_P}\nu\int_0^1e^{2Rv}Q(v)^2dv
 +\nu B_P\int_0^1e^{2Rv}\bigl(RQ(v)+Q'(v)\bigr)^2dv .
\]
Each integral $\int_0^1e^{2Rv}F(v)dv$ of a rational polynomial equals
$e^{2R}a(F)+b(F)$ with rational $a,b$, by the recursion
$\int_0^1e^{tv}v^kdv=e^ta_k+b_k$, $a_k=1/t-(k/t)a_{k-1}$,
$b_k=-(k/t)b_{k-1}$, $a_0=1/t$, $b_0=-1/t$. Hence $c=Ae^{2R}+B$ with exact
rationals $A,B$, and $\kappa=1-R^{-1}\log c$ needs one exponential and one
logarithm in ball arithmetic (Arb, $4000$ bits). The admissibility identities
are verified in exact rational arithmetic.

\begin{table}[ht]
\centering
\caption{Certified detector constants. Lower bounds are rounded down.}
\label{tab:kappa}
\begin{tabular}{llllll}
\toprule
$\nu$ & $R$ & $r$ & $c(P,Q,R,\nu)$ & $\kappa>$ & $\kappa/\nu>$\\
\midrule
$0.0199$ & $48.391$ & $0.7702$ & $5.1998\cdot10^{20}$ & $0.0142729911$ & $0.717235$\\
$0.0299$ & $32.244$ & $0.7708$ & $5.0471\cdot10^{13}$ & $0.0214483010$ & $0.717334$\\
$0.0339$ & $28.439$ & $0.7708$ & $1.1235\cdot10^{12}$ & $0.0243176419$ & $0.717334$\\
$0.0349$ & $27.625$ & $0.7709$ & $4.9777\cdot10^{11}$ & $0.0250349766$ & $0.717334$\\
$0.0399$ & $24.163$ & $0.7709$ & $1.5614\cdot10^{10}$ & $0.0286216496$ & $0.717334$\\
$0.0458$ & $21.05$ & $0.7708$ & $6.9431\cdot10^{8}$ & $0.0328539236$ & $0.717334$\\
$0.0499$ & $19.321$ & $0.7709$ & $1.2321\cdot10^{8}$ & $0.0357949954$ & $0.717334$\\
$0.0999$ & $9.648$ & $0.7707$ & $7.7588\cdot10^{3}$ & $0.0716640051$ & $0.717357$\\
\bottomrule
\end{tabular}
\end{table}

In our numerical optimization the ratio $\kappa/\nu$ did not exceed about
$0.7173$ for this mollifier shape at small $\nu$; the polynomials of Conrey,
Farmer, Kwan, Lin and Turnage-Butterbaugh~\cite{CFKL} give about $0.682\nu$.
Linear $Q$ does not suffice at these lengths: $Q=1-x$, $P=x$ needs
$\nu>0.1928$ for $\kappa>0$. The same code reproduces Young's
$c=2.35006777\ldots$ for $P=x$, $Q=1-x$, $R=1.3$, $\nu=1/2$, and Conrey's
$\kappa>0.4088$ for his polynomial of degree seven with $\nu=4/7$.

An independent program evaluates $Q$ by the Chebyshev recurrence directly
from the $x_j$ and $P$ by its series, separates the double integral by
Fubini into one-dimensional integrals, computes each by validated Arb
quadrature, and recomputes $c(\theta)$ and $h$ with $100$-digit interval
arithmetic in \texttt{mpmath}. Every enclosure overlaps the primary one, and
all thresholds hold again. With $\nu=\theta-1/2-10^{-4}$, the directed values
of $h(\theta;\kappa)$ at $\theta=0.534,0.535,0.54,0.5459,0.55$ exceed
$1.4806994\cdot10^{-5}$, $0.0015246940$, $0.0090231376$, $0.0177638490$ and
$0.0237708528$; at $\theta=0.5459$ the ratio to the rank-six Selberg value
$1.7764518\cdot10^{-5}$ exceeds $999.96$.

Numerical controls support Lemma~\ref{lem:count}. At $T=10^4,10^5,10^6$ and
for $Q$ of degrees $1,3,5,7$, including Conrey's polynomial and polynomials
with $\beta\ne1$, the direct evaluation of $\chi(s)V(1-s)$ agrees with
Lemma~\ref{lem:identity}, $\omega E_\beta$ is real and \eqref{eq:real} holds to
the working precision. On windows at these heights the sign changes of $Z$,
the level crossings of $\arg W_0$ and the zero counts coincide. For the test
functions $f=\zeta\prod_j((s-1/2)^2+\gamma_j^2)^{e_j}$ with planted double
($e_j=1$) and triple ($e_j=2$) critical zeros, the inequality of
Lemma~\ref{lem:count} holds with equality in the double-zero cases, double
zeros are not counted, and the variant that omits the on-line zeros of $\Vt$
fails for $\deg Q=1$.

\section{Scope}
The moment uses the attributed trilinear theorem of Bettin and Chandee,
classical special-function identities and Young's arithmetic evaluation.
The counting proof is localized with full weight for detector zeros on
the line. The simple-critical consequence additionally uses Wang's
recent short-interval pair theorem and the finite parity inequality.
Both analytic norm derivations retain their parameter losses and the
compact smoothing margin. Exact and interval controls verify finite
normalizations and constants; they do not substitute for the universal
analytical proof, establish external peer review or prove worldwide novelty.
All statements are asymptotic and give no effective height. Very large
fixed detector coefficients and smoothing orders can make constants
astronomical. The range does not reach positivity for every
$\theta>1/2$, and no universal simplicity statement or RH proof follows.

At $\theta=5339/10000$, $\nu=349/10000$, choose
$\eta=1/100000$. The narrower Gaussian exponent is $53389/100000$;
the two error exponents are exactly $-7/250000$ and $-937/400000$.
The two independent arithmetic paths give
$h>0.0003985233159135$. An exact rational check of 7,430 dyadic
configurations, together with the universal slope argument, detects
overlarge smoothing loss, zero margin, doubled separation loss,
uncontrolled tail and a reversed source parameter factor.

\medskip
\noindent\textbf{Data availability.} The programs, exact parameters, outputs
and independent replays are archived as computations EXP-010 and EXP-024--028 in the
repository~\cite{Record,Record28}.

\medskip
\noindent\textbf{Declaration of generative AI use.} Generative AI tools were
used for source research, derivation, review and implementation. Two analytical derivations and the stated numerical controls
were used to test the claims. This preprint has not undergone external peer review.

\begin{thebibliography}{99}
\bibitem{BC} S. Bettin and V. Chandee,
\emph{Trilinear forms with Kloosterman fractions}, Adv. Math.
\textbf{328} (2018), 1234--1262; arXiv:1502.00769v1, Theorem 1.
\bibitem{DLMF} NIST Digital Library of Mathematical Functions,
\url{https://dlmf.nist.gov/}, formulas 25.11.9, 13.4.4, 13.6.10, 10.27.8,
10.22.43 and 10.43.19, accessed 4 October 2026.

\raggedright

\bibitem{BCY}
H. M. Bui, J. B. Conrey and M. P. Young,
\emph{More than 41\% of the zeros of the zeta function are on the critical
line}, Acta Arith. \textbf{150} (2011), 35--64; arXiv:1002.4127.

\bibitem{CAOSShort}
F. Santib\'a\~nez-Leal,
\emph{Simple critical zeros in short intervals: stability, parity,
localization, Hilbert compression, and spectral defect},
Zenodo preprint (2026),
\url{https://doi.org/10.5281/zenodo.22860012}.

\bibitem{CFKL}
J. B. Conrey, D. W. Farmer, C.-H. Kwan, Y. Lin and
C. L. Turnage-Butterbaugh,
\emph{Short mollifiers of the Riemann zeta-function},
arXiv:2508.11108v1 (2025).

\bibitem{Conrey}
J. B. Conrey,
\emph{More than two fifths of the zeros of the Riemann zeta function are on
the critical line}, J. Reine Angew. Math. \textbf{399} (1989), 1--26.

\bibitem{Karatsuba}
A. A. Karatsuba,
\emph{On the zeros of the function $\zeta(s)$ on short intervals of the
critical line}, Math. USSR-Izv. \textbf{24} (1985), 523--537.

\bibitem{PearceCrump}
A. Pearce-Crump,
\emph{Optimising Selberg's method for critical zeros},
arXiv:2609.15329v1 (2026),
\url{https://arxiv.org/abs/2609.15329v1}.

\bibitem{Record}
F. Santib\'a\~nez-Leal,
\emph{EXP-010: localized Levinson parity transfer},
CAOS Research computational record (2026),
\url{https://github.com/fsantibanezleal/CAOS_RESEARCH/tree/main/problems/number-theory/riemann-hypothesis/experiments/EXP-010-levinson-parity-transfer}.

\bibitem{Record28}
F. Santib\'a\~nez-Leal,
\emph{EXP-028: separated short-window moment and independent arithmetic},
CAOS Research computational record (2026),
\url{https://github.com/fsantibanezleal/CAOS_RESEARCH/tree/main/problems/number-theory/riemann-hypothesis/experiments/EXP-028-chirp-separated-moment}.

\bibitem{Steuding}
J. Steuding,
\emph{On simple zeros of the Riemann zeta-function in short intervals},
Acta Math. Hungar. \textbf{96} (2002), 259--308.

\bibitem{SteudingThesis}
J. Steuding,
\emph{On simple zeros of the Riemann zeta-function},
doctoral dissertation, Universit\"at Hannover (1999).

\bibitem{WangShort}
B. Wang,
\emph{Simple critical zeros and distinct zeros of the Riemann zeta-function
in short intervals},
arXiv:2609.07918v1 (2026),
\url{https://arxiv.org/abs/2609.07918v1}.

\bibitem{Young}
M. P. Young,
\emph{A short proof of Levinson's theorem},
Arch. Math. \textbf{95} (2010), 539--548; arXiv:1002.4403.

\end{thebibliography}
\end{document}
